\chapter{القرآن كتاب هداية والاعتصام بالسنة}

\section{القرآن كتاب هداية}
القرآن تبيان لكل شيء، ولكن تبيان في ماذا؟ في أمر الهداية؛ فهو كتاب هداية، وليس كتاب جبر وهندسة وكيمياء وفيزياء. حتى لا يُقال: نحن نريد حلاً للمسألة الفلانية، فنقول: ﴿فَاسْأَلُوا أَهْلَ الذِّكْرِ إِن كُنتُمْ لَا تَعْلَمُونَ﴾، فيُسأل أهلها. وهذا أيضاً من الأصول العظيمة في القرآن. وقال تعالى: ﴿وَكَذَٰلِكَ أَوْحَيْنَا إِلَيْكَ رُوحًا مِّنْ أَمْرِنَا مَا كُنتَ تَدْرِي مَا الْكِتَابُ وَلَا الْإِيمَانُ وَلَٰكِن جَعَلْنَاهُ نُورًا نَّهْدِي بِهِ مَن نَّشَاءُ مِنْ عِبَادِنَا ۚ وَإِنَّكَ لَتَهْدِي إِلَىٰ صِرَاطٍ مُّسْتَقِيمٍ﴾.

\section{الوصية بالتمسك بالسنة}
فيا طالب الصراط المستقيم، إذا كنت تريد الصراط المستقيم، فعليك بسنة النبي الكريم محمد صلى الله عليه وسلم؛ استمسك بها، وعضّ عليها بالنواجذ، وادعُ إليها، واحذر من خلافها، تنل وتفز بجنة ربك وبرحمته.

نقف عند باب "كيف البيان". أسأل الله بأسمائه الحسنى وصفاته العلى أن ينفعني وإياكم بما قلنا وبما سمعنا.